
\begin{frame}
 
\frametitle{\resizebox{.95\linewidth}{!}{%
Addressing Dirichlet's Limitations: Extensions \& Alternatives}}
\note{Now, I want to talk about three useful versions of the Dirichlet distribution:

First, the Zero-Inflated Dirichlet. This model is good when we have many zeros in the data. The standard Dirichlet cannot handle zeros well, but this one can.

Second, the Nested Dirichlet. The standard Dirichlet assumes all parts are negatively related. This is not always true in real data. The nested version helps fix that.

Third, the Kummer–Dirichlet. This model is better when the data has outliers. It is more flexible and gives more stable results.}
\begin{itemize}
  \item[\textcolor{blue}{\large$\blacktriangleright$}] \textbf{ Zero-Inflated Dirichlet Models:}  
  Handle compositional data with many zeros.

  \item[\textcolor{blue}{\large$\blacktriangleright$}] \textbf{ Nested–Dirichlet Models:}  
  Address the limitation of standard Dirichlet where all components are negatively correlated.

  \item[\textcolor{blue}{\large$\blacktriangleright$}] \textbf{Kummer–Dirichlet Models:}  
  Offer greater robustness to outliers compared to the standard Dirichlet.
\end{itemize}


\end{frame}

\begin{frame}
\frametitle{Distributional Dirichlet Random Forests (DRF)}
\note{
Another idea that I really like is to fit an empirical probability function.

Actually, we can fit a probability function on the training dataset. Then, for a new observation, we can evaluate this probability function to obtain its predicted distribution.

Actually, we can fit a probability function on the training dataset. This method was introduced in 2022. I will not go into the details of how it works. I have a simple example to explain it, but I will skip it to save time.
}

\begin{equation*}
P(Y = y_i \mid X = \mathbf{x}_{\text{new}}) = w_i(\mathbf{x}_{\text{new}}), \quad i = 1, \dots, n, \quad \sum_{i=1}^{n} w_i(\mathbf{x}_{\text{new}}) = 1.
\end{equation*}

\vspace{0.5em}
\begin{block}{}
The formula for the normalized weight $w_i$ is:
\[
w_i(\mathbf{x}_{\text{new}})
= \frac{1}{B} \sum_{k=1}^B 
\frac{\mathbf{1}\!\left(\mathbf{x}_{\text{new}} \in L_k(i)\right)}{|L_k(i)|}.
\]
where $|L_k(i)|$ is the size of the terminal node in tree $k$ where $\mathbf{x}_{\text{new}}$ falls.
\end{block}
\end{frame}



\begin{frame}[plain]
  \begin{center}
    \vspace{2cm}

    {\Huge \textbf{Thank You!}\par}

    \vspace{1cm}

    \includegraphics[width=0.15\linewidth]{slides/DirichletRF.jpg}

    %\vspace{0.8cm}

    % QR code (requires: \usepackage{qrcode})
    \qrcode[height=1.8cm]{https://github.com/Xaleed/DirichletRF.git}

    \vspace{0.1cm}

    {\small QR code for the GitHub repository}

  \end{center}
\end{frame}
